\documentclass[11pt]{article}
\usepackage[margin=1in]{geometry}
\usepackage[T1]{fontenc}\usepackage[utf8]{inputenc}
\usepackage{lmodern,amsmath,amssymb,amsthm,mathtools,booktabs,microtype}
\usepackage[colorlinks=true,linkcolor=blue,citecolor=blue,urlcolor=blue]{hyperref}
\hypersetup{pdftitle={Agnostic oracle bounds for Gaussian neural networks with Lipschitz activations},pdfauthor={ChatGPT-6 Astra; Claude Opus 5.5; Kaarel Hänni; Dmitry Vaintrob}}
\newtheorem{theorem}{Theorem}[section]\newtheorem{lemma}[theorem]{Lemma}
\newtheorem{corollary}[theorem]{Corollary}\newtheorem{proposition}[theorem]{Proposition}
\theoremstyle{remark}\newtheorem{remark}[theorem]{Remark}
\newcommand{\E}{\mathbb E}\newcommand{\R}{\mathbb R}\newcommand{\Pp}{\mathbb P}
\newcommand{\KL}{\operatorname{KL}}\newcommand{\op}{\mathrm{op}}
\newcommand{\clip}{\operatorname{clip}}\newcommand{\dist}{\operatorname{dist}}
\newcommand{\whR}{\widehat R}\newcommand{\calF}{\mathcal F}
\newcommand{\calG}{\mathcal G}
\title{Agnostic oracle bounds for Gaussian neural networks\\with Lipschitz activations}
\author{ChatGPT-6 Astra \and Claude Opus 5.5 \and Kaarel H\"anni \and
Dmitry Vaintrob\thanks{Authors are listed in alphabetical order.}}
\date{17 November 2026}
\begin{document}\maketitle
\begin{abstract}
A small neural network may approximate a data distribution well even when no small network fits it exactly. We prove a coefficient-one agnostic oracle inequality for the ordinary posterior mean of a single dense Gaussian neural network. The activation can be any fixed globally Lipschitz function, including tanh, ReLU, GELU, sigmoid and softplus; it need not vanish at zero. Inputs have an arbitrary distribution on a cube, labels are bounded, and the likelihood is Gaussian. The guarantee holds at every finite learner width above comparator width and every Gaussian scale at or stronger than mean field. A data-dependent likelihood temperature is selected from a scalar free-energy equation, without knowing the comparator. For $H\ge2$ hidden layers, the excess-risk terms have rates $n^{-1/[2H(H+1)]}$ times a power of the comparator risk, and $n^{-1/(2H)}$ when that risk vanishes, with polynomial dependence on comparator width, coefficient bound and input dimension at fixed depth. The proof compresses typical Gaussian networks into finite dictionaries. It tracks low-rank feature coordinates together with a residual standard deviation, preserving a Lipschitz error recursion even for ReLU. These dictionaries are only proof devices: the learner retains one fixed dense architecture and independent Gaussian weights, with no prior over network sizes.
\end{abstract}

\section{What the guarantee says}
For squared loss, write $R(f)=\E(f(X)-Y)^2$. We seek a guarantee
\[
 R(\bar f_S)\le\inf_{g\in\calG_{m,B}}R(g)+\varepsilon
\]
after polynomially many observations, where $\calG_{m,B}$ consists of smaller networks of width $m$ and bounded coefficients. The learner itself may have arbitrarily larger width. The leading coefficient one matters: the claim preserves the approximation error achieved by the best comparator.

Our result applies to the actual finite Gaussian posterior, not only a zero-temperature or infinite-width limit. It covers arbitrary label noise and model misspecification. A coefficient bound is essential for this kind of distribution-uniform statement: width alone does not control the steepness of a real-valued network. The result is statistical and makes no claim that exact posterior inference or its free energy can be computed efficiently.

The proof has two parts. A neighborhood of a cloned comparator has sufficiently large Gaussian prior mass. Separately, with exponentially high prior probability, the learner can be approximated by one of finitely many functions whose number is controlled independently of learner width. The second step is the main issue near mean-field scaling: typical dense Gaussian matrices have Frobenius energy proportional to width, even though their operator norms stay bounded.

\section{Model and statement}
Fix input dimension $d\ge1$, hidden depth $H\ge1$, and total depth $L=H+1$. Let $\phi:\R\to\R$ be globally $A$-Lipschitz, where we take $A\ge1$, and put $\beta=|\phi(0)|$. All constants denoted $C$ or $c$ below may depend on $H,A,\beta$; different occurrences need not denote the same constant. Their dependence on $d,m,B,n,N$ is displayed or bounded polynomially.

Inputs lie in $[-1,1]^d$, labels in $[-1,1]$, and their joint Borel law $D$ is arbitrary. For $\bar x=(1,x)$ the width-$N$ learner is
\[
 h^1=\phi(W^1\bar x),\qquad
 h^\ell=\phi(W^\ell h^{\ell-1})\ (2\le\ell\le H),\qquad
 z_w=v^\top h^H.
\]
The activation is coordinatewise. Define
\[
 Q_N(w)=\frac dN\|W^1\|_F^2+\sum_{\ell=2}^H\|W^\ell\|_F^2
                  +N\|v\|_2^2.
\]
The prior $\mu_{N,\sigma}$ has density proportional to
$e^{-Q_N/(2\sigma^2)}$. Its coordinate variances are $N\sigma^2/d$ in the first layer, $\sigma^2$ internally and $\sigma^2/N$ in the readout. We assume
\begin{equation}
 0<\sigma^2\le1/N,\qquad \tau^2=N\sigma^2,\qquad \omega=1/\sigma^2=N/\tau^2.
 \label{eq:scales}
\end{equation}
This includes every $\sigma^2=N^{-\gamma}$ with $\gamma\ge1$, at every finite $N$.

For an IID sample $S$ of size $n$, put $\whR_S(f)=n^{-1}\sum_i(f(X_i)-Y_i)^2$. With $\lambda>0$, use the posterior
\begin{equation}
 d\Pi_\lambda(w)=Z_\lambda^{-1}
 e^{-\whR_S(z_w)/(2\lambda\sigma^2)}\,d\mu_{N,\sigma}(w).
 \label{eq:posterior}
\end{equation}
Equivalently, the observation likelihood is Gaussian with variance $n\lambda\sigma^2$ and mean $z_w(X)$. Prediction is the raw posterior mean $\bar f_\lambda=\int z_w\,d\Pi_\lambda$. Define its normalized free energy
\begin{equation}
 F_\lambda=-2\lambda\sigma^2\log Z_\lambda.
 \label{eq:freeenergy}
\end{equation}

The comparator class $\calG_{m,B}$ uses the same activation and hidden depth, width $m$, first-layer biases, no other biases, and coefficients of absolute value at most $B\ge1$. Section~\ref{sec:teacher} explains how to include biases in every layer.

Set
\[
 q_H=\begin{cases}3,&H=1,\\4H-2,&H\ge2,\end{cases}
 \qquad \rho_H=(q_H+2)/2,\qquad v_H=L\rho_H.
\]
For $H\ge2$, these are $\rho_H=2H$ and $v_H=2H(H+1)$.
For $H=1$, they are $5/2$ and $5$.

The calibration uses only the sample, fixed learner and confidence level.
Let $M\ge1$ be a depth- and activation-dependent upper bound on the native prior's mean squared loss, valid for every sample; Lemma~\ref{lem:moments} supplies it. There are sufficiently large constants $C,c$ depending only on $H,A,\beta$ for which the following choices work:
\begin{align}
 b_n&=C(d+1)^c+4\log\!\bigl(C(d+1)^c n^3\bigr),\\
 \theta_n&=\frac{C(d+1)^c[\log(C(d+1)n/\delta)]^c}{n},
 \qquad U_0=2M+1.
 \label{eq:calibration-data}
\end{align}
Their deliberately conservative dimension factors absorb the dictionary and norm-tail constants.

If $\theta_n b_n^{v_H}U_0^{v_H-\rho_H}\le1$, select the unique positive solution of
\begin{equation}
 \lambda_S^{v_H}
 =\theta_n[2F_{\lambda_S}+b_n\lambda_S]^{v_H-\rho_H}.
 \label{eq:root}
\end{equation}
Otherwise use $\lambda_S=\omega M$. We call these the root and fallback branches. The fallback handles sample counts below the quantitative threshold; it does not use a different prior or architecture.

\begin{theorem}[An observable risk certificate]
\label{thm:certificate}
The calibration is measurable and well defined. On the root branch,
$1/n\le\lambda_S\le1/b_n$. With probability at least $1-\delta$, simultaneously for every learner width $N\ge1$ and scale~\eqref{eq:scales},
\begin{equation}
 R(\bar f_{\lambda_S})
 \le\int R(z_w)\,d\Pi_{\lambda_S}(w)
 \le F_{\lambda_S}+B_d\lambda_S+n^{-3},
 \label{eq:certificate}
\end{equation}
where $B_d\le C(d+1)^c$ is a known constant from the proof.
On the fallback branch the deterministic certificate is $e^{1/2}M$.
\end{theorem}

\begin{theorem}[A risk-adaptive coefficient-one oracle]
\label{thm:oracle}
For each comparator width $m$ and cap $B$, let
\[
 R_*=\min_{g\in\calG_{m,B}}R(g).
\]
The explicit resource $\mathcal A_{m,B,n}$ in~\eqref{eq:fixed-resource} is polynomial in $d,m,B$ at fixed $H$, with only logarithmic dependence on $n$. Put
$\overline{\mathcal A}=\mathcal A_{m,B,n}+b_n+B_d+1+\log(4/\delta)$.
With probability at least $1-\delta$, simultaneously at every $N\ge m$ and every scale~\eqref{eq:scales}, the calibrated raw posterior mean satisfies
\begin{equation}
 R(\bar f_{\lambda_S})\le R_*+
 C\left\{\overline{\mathcal A}\,R_*^{\,1-1/L}\theta_n^{1/v_H}
       +\overline{\mathcal A}^{\,L}\theta_n^{1/\rho_H}\right\}
 +n^{-3}.
 \label{eq:oracle}
\end{equation}
The same bound holds for posterior-averaged raw risk.
The likelihood calibration does not use $m$, $B$ or $R_*$.
\end{theorem}

\begin{corollary}[Polynomial sample complexity]
At every fixed depth, a sample count polynomial in comparator width, coefficient bound, input dimension and inverse target excess suffices, uniformly over every finite $N\ge m$ and every $\gamma\ge1$. For $H\ge2$, the general-risk term decays as $n^{-1/[2H(H+1)]}$, and the zero-comparator-risk term as $n^{-1/(2H)}$, up to logarithms and the displayed polynomial resources. For one hidden layer the corresponding exponents are $1/5$ and $2/5$. These conservative exponents are not claimed optimal.
\end{corollary}

The posterior-averaged raw-risk bound also controls clipping the posterior mean, the average risk of clipped posterior draws, and the mean of clipped draws. The main prediction convention remains the ordinary raw mean.

A likelihood schedule fixed before seeing the data also gives a polynomial oracle with the same general-risk exponent; Appendix~\\ref{sec:fixed-calibration} proves this corollary. The evidence-based calibration improves adaptation to the comparator risk.

\section{A compression lemma that preserves a residual norm}
\label{sec:sampling}
The following construction avoids separately estimating a variance and then taking its square root. That separation would lose a square root of accuracy at each layer.

\begin{lemma}[Shared importance sampling]
\label{lem:sampling}
On a probability space let $b\in\R^r$ satisfy $\E bb^\top=I_r$, and let $Y(x)$ be a family of square-integrable variables. Suppose
\[
 \sup_x\E Y(x)^2\le M^2,\qquad
 q\ge1+\|b\|^2,\qquad |Y(x)|\le M\sqrt q,\qquad Z=\E q<\infty.
\]
Write
\[
 F(x)=\E bY(x),\qquad s(x)=\sqrt{\E Y(x)^2-\|F(x)\|^2}.
\]
For any probability measure $\xi$ on $x$, $0<\delta\le M/10$, and
\[
 k\ge C ZM^2\delta^{-2}\log(r+2),
\]
there are $k$ importance samples with the following property. Set
\[
 B_{j,*}=\sqrt{Z/k}\,b_j/\sqrt{q_j},\qquad
 y_j(x)=\sqrt{Z/k}\,Y_j(x)/\sqrt{q_j},\qquad
 G=B^\top B,\qquad \rho=(\delta/M)^2.
\]
Then the functions
\begin{align}
 \widetilde F&=(G+\rho I)^{-1/2}B^\top y,\qquad
 a=(G+\rho I)^{-1}B^\top y,\\
 \widetilde s&=\|(y-Ba,\sqrt\rho\,a)\|
 \label{eq:sampled-pair}
\end{align}
satisfy
\[
 \|(\widetilde F,\widetilde s)-(F,s)\|_{L^2(\xi)}\le C\delta.
\]
For every fixed choice of sample parameters, the map $y\mapsto(\widetilde F,\widetilde s)$ is $1$-Lipschitz.
\end{lemma}
\begin{proof}
Draw independently from the probability law $q\,dP/Z$. Put $R=Y-b^\top F$. Bessel's inequality gives $\|F(x)\|\le M$, $\E R(x)^2\le M^2$, and $\E bR=0$. Also $|R|/\sqrt q\le2M$.

Each summand $X=Zbb^\top/q$ is positive semidefinite, satisfies $X\preceq ZI$ and has mean $I$. Matrix Chernoff~\cite{tropp} gives, for $0<\kappa\le1$,
\[
 \Pp(\|G-I\|_\op>\kappa)\le2r\exp[-k\kappa^2/(3Z)].
\]
Let $e=B^\top(y-BF)$. Independence and $\|b\|^2/q\le1$ give
\[
 \E\|e\|_{L^2(\xi)}^2
 \le \frac Zk\E\int\frac{\|b\|^2R(x)^2}{q}\,d\xi(x)
 \le \frac{ZM^2}k.
\]
For a nonnegative bounded random variable $V$,
\[
 \E\left(\sqrt{k^{-1}\textstyle\sum_j V_j}-\sqrt{\E V}\right)^2
 \le\frac{\operatorname{Var}(V)}{k\E V}
 \le\frac{\|V\|_\infty}{k}.
\]
The case $\E V=0$ is immediate; otherwise rationalize the difference of square roots. Apply this pointwise to $V=ZR^2/q\le4ZM^2$ and integrate. It follows that
\[
 \E\big\|\|y-BF\|-s\big\|_{L^2(\xi)}^2\le4ZM^2/k.
\]
Matrix Chernoff and Markov's inequality leave positive probability that
$\|G-I\|_\op\le\delta/M$, $\|e\|_{L^2}\le C\delta$, and the last RMS error is at most $C\delta$. Fix such a sample. The spectrum of $G$ lies in $[1/2,3/2]$. Since $B^\top y=GF+e$, functional calculus gives $\|\widetilde F-F\|_{L^2}\le C\delta$.

The residual in~\eqref{eq:sampled-pair} is the distance of $(y,0)$ to the column space of $(B;\sqrt\rho I)$. At coefficient $F$ the residual vector is $(y-BF,-\sqrt\rho F)$. Its projection onto this column space has norm at most
$C(\|e\|+\rho M)$. Hence
\[
 \big|\widetilde s-\|y-BF\|\big|
 \le\sqrt\rho M+C\|e\|+C\rho M
\]
in $L^2(\xi)$, proving the claimed residual approximation.

Finally set $T=(G+\rho I)^{-1/2}B^\top$. Then $T^\top T\preceq I$ and direct expansion gives
$\widetilde s^2=\|y\|^2-\|Ty\|^2$. Thus the pair consists of $Ty$ and the norm of $(I-T^\top T)^{1/2}y$. Orthogonal energy decomposition and the reverse triangle inequality prove its contraction property.
\end{proof}

\begin{remark}[Additional empirical Gram constraints]
\label{rem:inputgram}
Suppose normalized sample vectors $u'$ satisfy $\|u'\|\le1$ and
$\E_{\rm imp}Zu'u'^\top\preceq I$. Matrix Chernoff also gives
\[
 \Pp\left(\left\|\frac Zk\sum_j u'_ju_j'^\top\right\|_\op>2\right)
 \le(\dim u')e^{-ck/Z}.
\]
Replacing the logarithm in Lemma~\ref{lem:sampling} by
$\log(r+\dim u'+2)$ permits this constraint simultaneously. The same applies to a bounded scalar second moment. These constraints will keep the sampled programs Lipschitz in their incoming feature coordinates.
\end{remark}

A simpler scalar version will be used at the last layer. If $\E c^2\le K$, choose
$q=1+\|a\|^2/K+c^2/K+G_0^2$, where $G_0$ is standard normal. Importance sampling of $\E cY$ has Hilbert-space variance at most
\[
 \frac Zk\,\E\frac{c^2Y^2}{q}\le\frac{ZK}k\,\E Y^2.
 \label{eq:scalar-variance}
\]
Thus $k\ge CZK\sup_x\E Y(x)^2/\varepsilon^2$ samples achieve $L^2(\xi)$ error at most $\varepsilon$, while the empirical Gram constraints in Remark~\ref{rem:inputgram} cost only logarithms.

\section{Finite dictionaries for typical Gaussian networks}
\label{sec:dictionary}
\begin{theorem}[A width-independent finite dictionary]
\label{thm:dictionary}
For $K\ge C(d+H+1)$ and $0<e<1/2$, there is a finite set $\calF_{K,e}$ of functions bounded by $M=C K^{L/2}$ such that
\begin{equation}
 \log|\calF_{K,e}|
 \le C(d+1)e^{-q_H}\log^c\!\bigl(C(d+1)K/e\bigr).
 \label{eq:dictionary-size}
\end{equation}
It depends on $d,H,\phi,K,e$, but not on learner width, prior scale or the measure below. For every probability measure $\xi$ on the input cube, every $N\ge1$ and every scale~\eqref{eq:scales},
\begin{equation}
 \mu_{N,\sigma}\left(\{\|z_w\|_\infty>M\}\cup\left\{
   \inf_{f\in\calF_{K,e}}\|z_w-f\|_{L^2(\xi)}>Me
   \right\}\right)
 \le C e^{-c\omega K}.
 \label{eq:dictionary-prob}
\end{equation}
The exceptional set may depend on $\xi$. The dictionary does not.
\end{theorem}

\subsection{Norm bounds and spectral coordinates}
The event
\begin{equation}
 \frac dN\|W^1\|_F^2\le K,\qquad
 \|W^\ell\|_\op^2\le K\ (2\le\ell\le H),\qquad
 N\|v\|^2\le K
 \label{eq:normevent}
\end{equation}
has probability at least $1-Ce^{-c\omega K}$. For the first and last blocks this is a chi-square bound. For an internal matrix, Gaussian operator-norm concentration around $2\tau$ gives the stated exponent once $K$ exceeds a fixed constant. The first-layer mean energy is $\tau^2(d+1)$, explaining the requirement $K\ge C(d+H+1)$.

On~\eqref{eq:normevent},
\[
 \|h^\ell(x)\|/\sqrt N\le U_\ell,\qquad
 U_1=A\sqrt{K(d+1)/d}+\beta,\qquad
 U_\ell=A\sqrt K\,U_{\ell-1}+\beta.
\]
Hence $U_\ell\le C K^{\ell/2}$ and $|z_w(x)|\le M=C K^{L/2}$.

Fix $\xi$. The normalized hidden kernel
$Q_\ell(x,y)=N^{-1}\langle h^\ell(x),h^\ell(y)\rangle$
is a positive finite-rank operator on $L^2(\xi)$, with trace at most $U_\ell^2$. At a cutoff $\eta_\ell>0$ retain its eigenvalues at least $\eta_\ell$. There is an $N$ by $r_\ell$ matrix $V_\ell$ with orthonormal columns such that
\begin{equation}
 F_\ell(x)=V_\ell^\top h^\ell(x)/\sqrt N,\qquad
 s_\ell(x)=\|(I-V_\ell V_\ell^\top)h^\ell(x)\|/\sqrt N,\qquad
 r_\ell\le U_\ell^2/\eta_\ell.
 \label{eq:state}
\end{equation}
The discarded covariance operator has norm at most $\eta_\ell$.
The pair $(F_\ell,s_\ell)$ has pointwise Euclidean norm at most $U_\ell$.
A Borel choice of spectral bases can be made for the finite Gram matrix, including repeated eigenvalues.

Conditional on earlier layers, the next row preactivation is exactly
\begin{equation}
 a_i^\top F_{\ell-1}(x)+Z_i(x),
 \label{eq:decomposition}
\end{equation}
where the $a_i$ are independent $N(0,\tau^2I)$ and the $Z_i$ are independent centered Gaussian fields, independent of the $a_i$, with covariance $\tau^2$ times the discarded kernel. Their pointwise variance is $\tau^2s_{\ell-1}(x)^2$.
On~\eqref{eq:normevent},
\[
 \frac1N\sum_i a_ia_i^\top\preceq KI,
\]
because the coefficient matrix is $\sqrt N W^\ell V_{\ell-1}$.

\subsection{Removing the Gaussian residual fields}
The next two bounds hold conditionally on all earlier layers and all $a_i$. Write
\[
 h_i(x)=\phi(a_i^\top F(x)+Z_i(x)),\quad
 m_i(x)=\E_Zh_i(x),\quad r_i(x)=h_i(x)-m_i(x).
\]
Suppose the discarded kernel has operator norm at most $\eta$ and pointwise diagonal at most $U^2$.

\begin{lemma}[Residual operator]
\label{lem:residual}
As an operator from $\R^N$ to $L^2(\xi)$, let
$B_rc=N^{-1/2}\sum_i c_i r_i$. Except on conditional probability $e^{-\omega K}$,
\begin{equation}
 \|B_r\|_\op\le C A\left\{\frac U{\sqrt\omega}
                      +\sqrt{\eta(K+1)}\right\}.
 \label{eq:residualop}
\end{equation}
Consequently every orthonormal $N$ by $r$ matrix $V$, even one chosen after the current features are seen, satisfies
\[
 \|V^\top(h-m)/\sqrt N\|_{L^2(\xi;\R^r)}
 \le\sqrt r\,\|B_r\|_\op.
\]
\end{lemma}
\begin{proof}
For fixed $\|c\|=1$, independence and Gaussian Poincar\'e give
$\E\|B_rc\|^2\le A^2\tau^2U^2/N=A^2U^2/\omega$.
As a function of the independent standard Gaussian row variables,
$\|B_rc\|$ is $A\tau\sqrt{\eta/N}$-Lipschitz. Indeed apply pointwise Cauchy--Schwarz in the row index, the Lipschitz property of $\phi$, and the operator norm $\sqrt\eta$ of the residual Gaussian covariance square root into $L^2(\xi)$.

Gaussian concentration~\cite{gaussian}, a $1/2$-net of the unit sphere of cardinality at most $5^N$, and the usual net bound on operator norm prove~\eqref{eq:residualop}. The squared deviation used in the union bound is
$2A^2\eta(K+\tau^2\log5)$: its tail exponent is $\omega K+N\log5$.
The last assertion follows by summing the squared bound over the columns of $V$.
\end{proof}

\begin{lemma}[A uniform residual standard deviation]
\label{lem:rms}
Let $P$ be an orthogonal projector of rank $r$ and put
\[
 T_P(x)=\|(I-P)h(x)\|/\sqrt N,\qquad
 b_P(x)=\sqrt{\|(I-P)m(x)\|^2/N+
                      N^{-1}\sum_i\operatorname{Var}(h_i(x))}.
\]
On the norm event, except on conditional probability at most $Ce^{-\omega K}$, simultaneously for every such projector of every allowed rank,
\begin{align}
 \|T_P-b_P\|_{L^2(\xi)}
 \le C A\biggl\{
 U\sqrt{\frac{r+1}{\omega}}
 +\sqrt{\eta\left[K+(r+1)
                  \log\left(2+\frac{\sqrt{r+1}}{\kappa}\right)\right]}
 +\kappa U_{\rm next}
 \biggr\}
 \label{eq:rms}
\end{align}
for $0<\kappa<1$. Here $U_{\rm next}$ is the next-layer norm bound, and constant factors in it are harmless.
\end{lemma}
\begin{proof}
First fix $P$ independently of the current residual fields. The map from those standard Gaussian variables to $\|T_P-b_P\|_{L^2(\xi)}$ is $A\tau\sqrt{\eta/N}$-Lipschitz. Pointwise Poincar\'e gives
\[
 \operatorname{Var}(T_P(x))
       \le A^2\tau^2s_{\rm prev}(x)^2/N.
\]
The row fluctuations are independent, so
\[
 \E T_P(x)^2=b_P(x)^2-
                 N^{-1}\sum_iP_{ii}\operatorname{Var}(h_i(x)).
\]
The last term is between zero and $rA^2\tau^2s_{\rm prev}(x)^2/N$.
Using $\E(T-\sqrt{\E T^2})^2\le2\operatorname{Var}(T)$ and the square-root inequality gives
\[
 \E\|T_P-b_P\|_{L^2}\le C A U\sqrt{(r+1)/\omega}.
\]

A $\kappa$-net of the rank-$r$ Stiefel manifold in Frobenius norm has cardinality at most $(1+2\sqrt r/\kappa)^{Nr}$. Its projectors approximate in operator norm to within $2\kappa$. On the norm event, $\|h\|/\sqrt N\le U_{\rm next}$; the coefficient Gram bound also gives $\|m\|/\sqrt N\le C U_{\rm next}$. Replacing $P$ by a net projector therefore changes both $T_P$ and $b_P$ in $L^2$ by at most $C\kappa U_{\rm next}$. Gaussian concentration and this net prove~\eqref{eq:rms}. The net contribution to the squared deviation is at most $C\eta\tau^2r\log(2+\sqrt r/\kappa)$, bounded by the displayed expression.

For simultaneous ranks, assign failure probability $e^{-\omega K-r-1}$ at rank $r$ and sum. The additional term in the squared deviation is bounded by $C\eta(r+1)$ and is already covered. Ranks cannot exceed $N$. The norm-event intersection handles the interpolation bounds; its complement was counted separately in~\eqref{eq:normevent}.
\end{proof}

For the actual spectral projector at layer $\ell$, set
$b_i=\sqrt N(V_\ell)_{i,*}$, so $N^{-1}\sum_i b_ib_i^\top=I$.
Let $G_0$ be an independent scalar standard normal and define
\[
 Y_i(x,G_0)=\phi(a_i^\top F_{\ell-1}(x)+\tau s_{\ell-1}(x)G_0).
\]
The two lemmas approximate the actual pair $(F_\ell,s_\ell)$ by
\begin{equation}
 \left(\E_{i,G_0}b_iY_i,\quad
       \sqrt{\E_{i,G_0}Y_i^2-\|\E_{i,G_0}b_iY_i\|^2}\right).
 \label{eq:modelpair}
\end{equation}
Although $b_i$ depends on the current residual fields, the factor bound was uniform over orthonormal matrices and the RMS bound uniform over projectors. Thus this substitution does not assume false conditional independence.

\subsection{Why errors do not acquire repeated square roots}
The coordinate functions of $b_i$ are orthonormal in
$L^2(\text{uniform }i\times\text{Gaussian }G_0)$. For any scalar $Y$, the map
\[
 Y\longmapsto\left(\E bY,\ \|Y-\operatorname{proj}_bY\|_2\right)
\]
is $1$-Lipschitz by Pythagoras and the reverse triangle inequality.
Together with the coefficient Gram bound, changing a previous state $(F,s)$ to $(F',s')$ changes~\eqref{eq:modelpair} by at most
\begin{equation}
 A\sqrt{K\|F-F'\|_{L^2}^2+\tau^2\|s-s'\|_{L^2}^2}.
 \label{eq:propagation}
\end{equation}
No derivative of $\phi$ is used. In particular the argument never differentiates a ReLU Gaussian mean with respect to variance at zero.

Apply Lemma~\ref{lem:sampling} to~\eqref{eq:modelpair}, evaluated on the \emph{already approximated} preceding state, with
\[
 q=1+\|a\|^2/K+\|b\|^2+G_0^2,\qquad
 Z\le2+r_{\rm prev}+r.
\]
Then $Y/\sqrt q=\psi_c(\sqrt K\,a'^\top F+g's)$, where
\[
 \psi_c(t)=c\phi(t/c),\quad c=1/\sqrt q,\quad
 a'=a/(\sqrt K\sqrt q),\quad g'=\tau G_0/\sqrt q .
\]
All normalized parameters have norm at most one; the normalized feature and its second moment are bounded by $C U_\ell$. The vector
$u'=(a',g')$ satisfies the empirical-Gram conditions of Remark~\ref{rem:inputgram}, since its importance-weighted expected Gram is
$\operatorname{diag}(\E aa^\top/K,\tau^2)\preceq I$.
Consequently the selected sampled layer, including Gram whitening and radial projection to radius $U_\ell$, is $C A\sqrt K$-Lipschitz in its incoming state, independently of its ranks.

At the first layer apply the same sampling lemma directly to the empirical distribution of actual incoming rows and $b_i$, with
$q=1+(d+1)\|W^1_{i,*}\|^2/K+\|b_i\|^2$.
There is no Gaussian residual-removal error at that layer.

At the final hidden layer a scalar approximation is enough. Write the readout as $N^{-1}\sum_i c_i h_i$ with $c_i=Nv_i$ and $N^{-1}\sum_i c_i^2\le K$. Lemma~\ref{lem:residual} gives
\[
 \left\|z-\E_{i,G_0}c_iY_i\right\|_{L^2}
 \le C\sqrt K\left\{U_{H-1}/\sqrt\omega+
                         \sqrt{\eta_{H-1}(K+1)}\right\}.
\]
The scalar importance construction after Lemma~\ref{lem:sampling} compresses this expectation. Its selected empirical input Gram and coefficient second moment ensure state Lipschitz constant $C A K$. For $H=1$ apply this scalar construction directly to the first-layer row distribution.

\subsection{Quantitative ranks and a finite dictionary}
We give the bookkeeping because the polynomial degree is part of the claim.
For $H\ge2$ set the final absolute output accuracy to $\varepsilon=C K^{L/2}e$.
Take local state errors
\[
 \delta_\ell=c_H\varepsilon K^{-(H-\ell+1)/2},\qquad
 \eta_{H-1}=c_H\varepsilon^2/K^2,
\]
and choose successively
\[
 \eta_{\ell-1}
 =c_H\frac{\delta_\ell^2\eta_\ell}{K U_\ell^2 J},
 \qquad J=C_H\log(C_HK/e).
\]
Small constants are chosen in depth order to absorb the Lipschitz factors in~\eqref{eq:propagation}. With $\kappa_\ell=c_H\delta_\ell/U_\ell$, the residual errors in Lemmas~\ref{lem:residual} and~\ref{lem:rms} are at most $\delta_\ell$, apart from their finite-effective-width terms.
The deterministic rank caps satisfy
\begin{equation}
 R_\ell\le C_H e^{-2(H-\ell)}J^{H-\ell-1},
                  \qquad 1\le\ell\le H-1.
 \label{eq:rankcaps}
\end{equation}
The importance sample counts can be chosen as
\[
 k_\ell\le C_H(1+R_{\ell-1}+R_\ell)e^{-2}\log^c(C_HK/e),
\]
with $Z_1\le3+R_1$ at the first layer and the corresponding scalar count at the terminal layer. The total scalar parameter count, including the normalized incoming and outgoing coefficient vectors, is bounded by
\begin{equation}
 C_H(d+1)e^{-(4H-2)}\log^c(C_HK/e).
 \label{eq:paramcount}
\end{equation}
For example the largest product is of order $e^{-2}R_1^2$; subsequent layers and the scalar terminal fit the same bound.

The finite-effective-width terms are small whenever
\begin{equation}
 \omega\ge T:=C_H K^{-1}e^{-2}
              \max(1,R_2,\ldots,R_{H-1})\log^c(C_HK/e).
 \label{eq:effectivewidth}
\end{equation}
The first layer contributes no such condition. If $\omega<T$, then $N\le\omega<T$, and we instead cover the original finite network directly on~\eqref{eq:normevent}. The number of parameters is at most
$C_H(T^2+dT)$, which fits~\eqref{eq:paramcount}. Entrywise parameter bounds and network Lipschitz constants are polynomial in $K,T,d$ at fixed depth; parameter quantization adds only logarithms. This small-width argument does not change the prior or introduce a mixture of widths. For $H=1$ the direct scalar sampling argument has no effective-width threshold.

It remains to discretize the scalar perspective functions without assumptions on the asymptotic behavior of $\phi$. Each $\psi_c$ is $A$-Lipschitz and $|\psi_c(0)|\le\beta$. On a bounded interval $[-T_0,T_0]$, the class of all such scalar functions has a uniform $\eta$-net with logarithmic cardinality
\[
 C(1+AT_0/\eta)\log(2+(\beta+AT_0)/\eta).
\]
Construct it by a mesh of spacing $\eta/(CA)$ and quantized vertex values, with piecewise-linear interpolation. The representatives can be taken $2A$-Lipschitz. This factor is absorbed into the depth-dependent constants.

At layer $\ell<H$ the argument interval has radius $T_0\le C U_\ell$.
Choose scalar activation accuracy
$\eta_\ell'=c_H\varepsilon K^{-(H-\ell+1)/2}/\sqrt{Z_\ell}$;
at the terminal layer choose
$\eta_H'=c_H\varepsilon/\sqrt{KZ_H}$.
The empirical input-Gram constraints and the state Lipschitz bounds above propagate these errors to total output error at most $C\varepsilon$. In both cases
\[
 T_0/\eta_\ell'\le C_H\sqrt{Z_\ell}/e.
\]
Thus independently choosing the scalar activation at every atom costs at most
\[
 \sum_\ell C_H k_\ell\sqrt{Z_\ell}/e
 \le C_H e^{-3H}\log^c(C_HK/e)
\]
in log cardinality. For $H\ge2$, $3H\le4H-2$, so~\eqref{eq:paramcount} still dominates.
For $H=1$, bounded $Z$ and $O(e^{-2})$ scalar samples give exponent three.

Finally quantize the remaining bounded program parameters. The ridge is bounded below by a polynomial in $e$ and $K^{-1}$, so matrix inversion, inverse square root and the augmented residual formula have polynomial parameter Lipschitz constants. State radial projections are contractions, and the chosen scalar functions have bounded Lipschitz constants. A sufficiently fine polynomial parameter mesh therefore costs only additional logarithms. Representatives can be selected from nonempty mesh cells intersecting the empirical-Gram constraint set, preserving its state Lipschitz bounds. Clip the terminal output to $[-M,M]$.

The resulting finite set is independent of $\xi,N,\sigma$. Conditional failure probabilities from the Gaussian lemmas sum over the fixed number of layers to $Ce^{-c\omega K}$. Intersecting with~\eqref{eq:normevent}, choosing small depth-dependent approximation constants, and the direct small-width cover prove Theorem~\ref{thm:dictionary}.

\section{A Gaussian neighborhood of a smaller network}
\label{sec:teacher}
This part uses only the Lipschitz property of the activation; no homogeneity or boundedness is required.
For a width-$m$ teacher with coefficients bounded by $B\ge1$, define
\begin{align}
 R_0&=1+\beta+A(d+m+1)B,\qquad U_{m,B}=mB R_0^H,\\
 E_{m,B}&=2\{d(d+1)+Hm^2\}B^2,\qquad D_{m,B}=d+2+(H-1)m.
 \label{eq:teacher-costs}
\end{align}
Choose $b_H=C\sqrt{\log(2H+2)}$ sufficiently large and put
\begin{align}
 T_0&=1+A(\sqrt2mB+b_H)+R_0,\\
 P_{m,B}&=(\sqrt2mB+1)A(d+1+H\sqrt m)T_0^{H-1}+R_0^H,\\
 \mathcal A_{m,B}(\lambda)
 &=E_{m,B}+4+
       2D_{m,B}\log\!\left(\frac{32P_{m,B}(U_{m,B}+1)}{\lambda}\right).
 \label{eq:Ag}
\end{align}
Enlarging a numerical constant in $P_{m,B}$ is harmless if needed.
These expressions have polynomial dependence on $d,m,B$ at every fixed $H$.

\begin{lemma}[Comparator neighborhood]
\label{lem:teacher}
For every $g\in\calG_{m,B}$, every $N\ge m$, every scale~\eqref{eq:scales} and $0<\lambda\le1$,
\begin{equation}
 F_\lambda\le\whR_S(g)+\lambda\mathcal A_{m,B}(\lambda).
 \label{eq:teacher-comparison}
\end{equation}
This holds deterministically for every sample.
\end{lemma}
\begin{proof}
Partition each width-$N$ hidden layer into $m$ groups of sizes
$k_j\in\{\lfloor N/m\rfloor,\lceil N/m\rceil\}$. Copy the teacher's first-layer rows, divide every internal edge leaving group $j$ by $k_j$, and divide each readout coefficient by its group size. This clone computes $g$ exactly for every activation. Since $N/k_j\le2m$, its energy satisfies $Q_N(w_*)\le E_{m,B}$. Its hidden coordinates at layer $\ell$ have absolute value at most $R_0^\ell$, and $|g|\le U_{m,B}$.

Use standard Gaussian coordinates in each parameter block. Let $P$ project onto the subspace of group-constant vectors, with orthonormal basis $u_j=k_j^{-1/2}\mathbf1_{\mathrm{group}\ j}$. For $0<\eta\le1/2$, require absolute value at most $\eta$ for every first-layer and readout standardized coordinate and every internal row coefficient along each $u_j$. There are $ND_{m,B}$ independent constraints. A standard normal assigns probability at least $3\eta/8$ to $[-\eta,\eta]$, so their joint mass is at least $(3\eta/8)^{ND_{m,B}}$.

Conditioning a standard normal to such a symmetric interval preserves the unit-subgaussian bound: by symmetry its exponential moment is an expectation of $\cosh(s|Z|)$, and the conditioned $|Z|$ is stochastically no larger than the original one. Orthogonal unconstrained coordinates remain Gaussian and independent. Thus every centered linear form in the standardized internal matrix is unit-subgaussian after conditioning. Two $1/4$ sphere nets imply
\[
 \Pp(\|Z\|_\op>2u\sqrt N\mid\text{coordinate constraints})
 \le2\exp[(2\log9-u^2/2)N].
\]
The choice of $b_H$ therefore permits intersecting with
$\|Z\|_\op\le b_H\sqrt N$ at every internal layer while retaining at least half the mass. Call the resulting centrally symmetric event $A_\eta$.

Translate $A_\eta$ by the standardized clone. Perturbations in the first-layer and readout coordinates are bounded by $\eta/\sqrt d$ and $\eta/N$, respectively. For an internal perturbation $\Delta W=\tau Z/\sqrt N$,
\[
 \|\Delta W\|_\op\le b_H,\qquad
 \|\Delta WP\|_\op\le\eta\sqrt m.
\]
The part on $I-P$ annihilates the clone's hidden vector. Normalized hidden errors $e_\ell=\|h^\ell-h_*^\ell\|/\sqrt N$ consequently obey
\[
 e_1\le A\eta(d+1),\qquad
 e_\ell\le A(\sqrt2mB+b_H)e_{\ell-1}
                         +A\eta\sqrt m\,R_0^{\ell-1}.
\]
The normalized clone readout norm is at most $\sqrt2mB$; the perturbation readout has $\ell^1$ norm at most $\eta$. The output error is therefore at most
\[
 (\sqrt2mB+1)e_H+\eta R_0^H\le\eta P_{m,B}.
\]

For a standard Gaussian measure $G$ and a centrally symmetric event $A$,
\[
 G(A+a)\ge G(A)e^{-\|a\|^2/2}.
\]
Indeed the density ratio on $A$ is
$e^{-\|a\|^2/2}e^{-\langle Z,a\rangle}$ and its conditional expectation is at least one by Jensen and symmetry. Here
$\|a\|^2=\omega Q_N(w_*)\le\omega E_{m,B}$.
Since $N\le\omega$ and $\omega\ge1$, the prior mass of the uniform-output neighborhood of radius $\rho=\eta P_{m,B}$ is at least
\[
 \exp\!\left[-\omega\left\{
   E_{m,B}/2+D_{m,B}\log(8/(3\eta))+\log2\right\}\right].
\]
Choose $\rho=\lambda/[4(U_{m,B}+1)]$ and $\eta=\rho/P_{m,B}$, decreasing $\eta$ by a harmless numerical factor if necessary. Conditioning the prior on this neighborhood gives a comparison law $\nu$ with
\[
 \int\whR_S(z_w)\,d\nu\le\whR_S(g)+\lambda,\qquad
 2\sigma^2\KL(\nu\Vert\mu)\le\mathcal A_{m,B}(\lambda)-1.
\]
The Gibbs variational identity
\[
 F_\lambda=\inf_\nu\left\{\int\whR_S(z_w)\,d\nu+
                                  2\lambda\sigma^2\KL(\nu\Vert\mu)\right\}
\]
proves~\eqref{eq:teacher-comparison}.
\end{proof}

\begin{corollary}[One resource for every allowed calibration]
\label{cor:fixedteacher}
For $n\ge2$, define
\begin{equation}
 \mathcal A_{m,B,n}
 =E_{m,B}+4+
  2D_{m,B}\log\!\left(32nP_{m,B}(U_{m,B}+1)\right).
 \label{eq:fixed-resource}
\end{equation}
Then for every sample and every $\lambda\ge1/n$,
\[
 F_\lambda\le\whR_S(g)+\lambda\mathcal A_{m,B,n}.
\]
Moreover $U_{m,B}^2\le C_{H,A,\beta}E_{m,B}^{L}$.
\end{corollary}
\begin{proof}
Use the preceding proof with the fixed neighborhood radius
$1/[4n(U_{m,B}+1)]$. Its empirical-risk increase is at most $1/n\le\lambda$, while its normalized entropy cost is independent of $\lambda$. For the last bound,
$R_0\le C(d+m+1)B$ and
$E_{m,B}\ge c(d+m+1)^2B^2$, with activation- and depth-dependent constants. Thus
$m^2B^2R_0^{2H}\le C E_{m,B}^{H+1}$.
\end{proof}

\begin{remark}[Biases in every layer]
If $\phi$ is not identically zero, choose a scalar $t_*$ with
$c_*=\phi(t_*)\ne0$. An extra neuron can carry the constant $c_*$ through the network: use first-layer bias $t_*$ and subsequent copying edge $t_*/c_*$. Replace each internal bias $b$ by an edge of weight $b/c_*$ from this neuron, and do the same for a readout bias. Thus ordinary biased comparators are included with width $m+1$ and coefficient cap at most $C_\phi B$. The learner and its prior are unchanged. If $\phi$ is identically zero, every network in the stated architecture outputs zero and the theorem is immediate.
\end{remark}

\section{Native moments and calibrated posterior risk}
\begin{lemma}[Uniform native moments]
\label{lem:moments}
There is $C=C_{H,A,\beta}$ such that, for every input and every scale~\eqref{eq:scales},
\[
 \E_\mu |z_w(x)|^4\le C,\qquad
 \E_\mu R(z_w)^2\le C,\qquad
 \E_\mu\whR_S(z_w)\le C.
\]
Consequently $Z_\lambda\ge e^{-C/(\lambda\sigma^2)}$ for every sample.
\end{lemma}
\begin{proof}
The first-layer preactivation variance is at most $2$. The inequality
$|\phi(t)|\le A|t|+\beta$, Gaussian fourth moments and
$(N^{-1}\sum_i h_i^2)^2\le N^{-1}\sum_i h_i^4$ inductively bound
$\E(N^{-1}\|h^\ell\|^2)^2$ by a depth-dependent constant. Conditional on the final hidden layer, the readout is Gaussian with variance
$\sigma^2\|h^H\|^2/N$. Since $\sigma^2\le1$, its fourth moment has the claimed bound. Jensen and the bounded labels give the risk moment statements. Jensen applied to the exponential defining $Z_\lambda$ gives its lower bound.
\end{proof}


\section{A localized statistical argument}
\label{sec:localized-direct-oracle}
We now prove Theorems~\ref{thm:certificate} and~\ref{thm:oracle}. The finite dictionary from Theorem~\ref{thm:dictionary}, expressed at absolute accuracy $0<a\le1$, satisfies
\[
 \log|\mathcal D(K,a)|
 \le B K^p a^{-q}\left[\log(C(d+1)K/a)\right]^s,\qquad
 q=q_H,\quad p=Lq/2,\quad B=C(d+1),
\]
after increasing the logarithmic exponent. Its good set $G_\xi(K,a)$ has prior complement at most $Ce^{-c\omega K}$, and both actual good outputs and dictionary members have amplitude at most $CK^{L/2}$. The initial level is $K_0=C(d+H+1)$.
Take $\rho=(q+2)/2$ and $v=p+L=L\rho$; since $L\ge2$, $v\ge2\rho$.
All constants arising from $K_0$ below are polynomial in $d$ at fixed depth and are absorbed in~\eqref{eq:calibration-data}. The proof constants $B_0,B_1,C_0,C_1$ are chosen with this polynomial dependence, and $B_0$ is the constant denoted $B_d$ in the theorem. We use the notation $F_S(\lambda)=F_\lambda$.

The sets can be chosen measurable: the actual networks are continuous in parameters and inputs, the spectral matrices have Borel eigenbases, and membership in a finite-dictionary $L^2$ neighborhood is measurable. In particular this remains true for $\xi=(D_X+\widehat D_{X,S})/2$. The good set may depend on the sample; only the dictionary used in the statistical union bound must be fixed.

\subsection{One sample event for all calibrations}
Take $U_0,b_n,\theta_n$ as in~\eqref{eq:calibration-data}, with constants enlarged to dominate the dictionary constants and confidence allocation. Then the following learner is well defined. If
$\theta_n b_n^v U_0^{v-\rho}\le1$, choose the positive solution of
\begin{equation}
 \lambda_S^v=\theta_n[2F_S(\lambda_S)+b_n\lambda_S]^{v-\rho}.
 \label{eq:localized-evidence-root}
\end{equation}
Otherwise choose $\lambda_S=\omega M$.
The root is unique and measurable because
$\lambda^v/[2F_S(\lambda)+b_n\lambda]^{v-\rho}
=\lambda^\rho/[2F_S(\lambda)/\lambda+b_n]^{v-\rho}$
is continuous and strictly increasing from zero to infinity.
Here $F_S\le M$ and $F_S(\lambda)/\lambda$ is nonincreasing.
On the root branch, enlarging the constants in~\eqref{eq:calibration-data} if necessary gives
$1/n\le\lambda_S\le1/b_n$ and
$u_S:=2F_S(\lambda_S)+b_n\lambda_S\le U_0$.

For use in the proof, prebuild dyadic norm levels $K_j=K_0 2^j$
through $C_0n$ and dyadic accuracies from $c_0/n$ through one.
There is a single sample event of probability at least $1-\delta/2$
with the following property. Whenever $\lambda\ge1/n$,
$0<u\le U_0$, $\lambda K\le C_1u$, and $a$ lies within a factor two
of $c_1\lambda K/\sqrt u$, each network good at $(K,a)$ with
$\widehat R_S(z)\le u$ satisfies
\begin{equation}
 R(z)\le\widehat R_S(z)+\kappa\lambda K
 \label{eq:localized-level-deviation}
\end{equation}
provided
\begin{equation}
 n\ge C B[\log(Cn/\delta)]^{s+4}
       u^\rho\lambda^{-2\rho}K^{v-2\rho}.
 \label{eq:localized-level-budget}
\end{equation}
Here $\kappa>0$ is a sufficiently small fixed constant, and $c_1$
is chosen after $\kappa,C_1,U_0$, so all required accuracies are at most
one and at least $c_0/n$.

Indeed, apply finite-class Bernstein to the squared losses of every
fixed dictionary, allocating confidence over the finite level/accuracy
list. A dictionary member has loss range $O(K^L)$ and variance at most
that range times its population risk. Inverting Bernstein gives
\[
 R(f)\le\widehat R(f)+C\sqrt{\widehat R(f)K^L t/n}+C K^L t/n,
 \quad t=\log|\mathcal D|+\log(1/\delta_{j,a}).
\]
The $L^2((D_X+\widehat D_X)/2)$ approximation controls both empirical
and population root risks. At the selected accuracy the approximant
has empirical risk $O(u)$, and transferring root risks costs
$O(\sqrt u\,a+a^2)=O(c_1\lambda K)$.
The leading Bernstein condition is exactly
\eqref{eq:localized-level-budget}; its linear term is absorbed using
$\lambda K\le C_1u$. Constants can be chosen to yield
\eqref{eq:localized-level-deviation}.
Maximizing the right side of \eqref{eq:localized-level-budget} at
$K\le C_1u/\lambda$, using $v\ge2\rho$, gives
$C B\log^{s+4}(Cn/\delta)u^{v-\rho}\lambda^{-v}$.
Thus the evidence root enforces all required level budgets.

\subsection{A sample-dependent complexity with uniform prior tails}
Fix a sample on the preceding statistical event and its selected root
$\lambda,u$. For every dyadic level $K_j$ whose value can enter
$K_j\le u/(\kappa\lambda)$, choose the prebuilt accuracy nearest below
$c_1\lambda K_j/\sqrt u$. At larger levels set the accuracy to one.
Use the good sets with $\xi_S=(D_X+\widehat D_X)/2$ and these accuracies.
Define
\[
 Q_S(w)=\max\{K_0,\ 2\sup\{K_j:w\notin G_{\xi_S}(K_j,a_j)\}\},
\]
where the supremum of the empty set is zero. For each fixed sample the
prior tail bound and a countable union give
\[
 \mu(Q_S>t)\le C'e^{-c'\omega t}\quad(t\ge2K_0).
\]
Only finitely many levels fail almost surely, so $Q_S$ is finite
almost surely. If $Q_S=K_j$, then $w\in G_{\xi_S}(K_j,a_j)$.
In particular $R(z_w)\le C(1+Q_S^L)$ almost surely.
Choose $\kappa$ small enough that tail integration gives, pointwise
for every sample,
\begin{equation}
 \log\mathbb E_\mu e^{\kappa\omega Q_S}\le\omega B_0,
 \qquad
 \mathbb E_\mu[R(z_w)e^{\kappa\omega Q_S/4}]
       \le C_0e^{\omega B_1}.
 \label{eq:localized-complexity-mgf}
\end{equation}
The constants do not depend on the sample, selected root, width, or
scale. Polynomial factors in $Q_S$ are absorbed by a smaller
exponential rate, using $\omega\ge1$.
Measurability follows from the measurable good-set tests and countable
supremum. No data-dependent complexity is placed inside an
unjustified empirical-process supremum; its prior inequalities are
pointwise in the sample.

Let $J=\widehat R_S+\kappa\lambda Q_S$. On $J>u$,
\[
 e^{-\omega\widehat R_S/(2\lambda)}
 \le e^{-\omega u/(4\lambda)}e^{\kappa\omega Q_S/4}.
\]
After normalization and \eqref{eq:localized-complexity-mgf}, the
posterior risk on this set is at most
$C_0\exp[\omega F_S/(2\lambda)-\omega u/(4\lambda)+\omega B_1]$,
which is at most $n^{-3}$ by the choice of $b_n$.
On its complement, \eqref{eq:localized-level-deviation} applies at
$K=Q_S$ and gives $R\le\widehat R_S+\kappa\lambda Q_S$.
Entropy duality in \eqref{eq:localized-complexity-mgf} gives
\[
 \kappa\lambda\mathbb E_\Pi Q_S
 \le(\lambda/\omega)\operatorname{KL}(\Pi\Vert\mu)+B_0\lambda.
\]
Integrate, extend the nonnegative right side to the entire posterior,
and use the variational identity. This proves the observable certificate
\begin{equation}
 \int R(z_w)d\Pi_{S,\lambda_S}
 \le F_S(\lambda_S)+B_0\lambda_S+n^{-3}.
 \label{eq:localized-certificate}
\end{equation}
The same holds for the ordinary raw posterior mean by Jensen. On the
fallback branch, Jensen gives $Z\ge e^{-1/2}$, so the deterministic
certificate $e^{1/2}M$ is valid instead.

\subsection{The resulting oracle powers}
Suppose a fixed teacher $g$ has a deterministic comparison resource
$A\ge1$ with
$F_S(\lambda)\le\widehat R_S(g)+\lambda A$ for every sample and
$\lambda\ge1/n$. Splitting the two terms in the root yields
\[
 \lambda_S\le C\left[
 \theta_n^{1/v}\widehat R_S(g)^{1-\rho/v}
 +\theta_n^{1/\rho}(2A+b_n)^{v/\rho-1}\right].
\]
Consequently, with $\bar A=A+b_n+B_0+1$, the empirical oracle is
\[
 \widehat R_S(g)+C\left[
 \bar A\widehat R_S(g)^{1-\rho/v}\theta_n^{1/v}
 +\bar A^{v/\rho}\theta_n^{1/\rho}\right]+n^{-3}.
\]
If $v=L\rho$ and $\|g\|_\infty^2\le C A^L$, Bernstein's scalar
teacher concentration is absorbed by those same terms, after including
confidence logarithms in $\bar A$. To check leading coefficient one,
put $D=\bar A^L\theta_n^{1/\rho}$. Bernstein gives
$\widehat R(g)-R(g)\le C\sqrt{R(g)D}+CD$, while
$\sqrt{RD}\le R^{1-1/L}D^{1/L}+D$ for $L\ge2$.
Also $\widehat R(g)\le2R(g)+CD$, which handles the sublinear power
without multiplying the leading population-risk term. Thus
\begin{equation}
 R(\mathbb E_\Pi z)\le R(g)+C\left[
 \bar A R(g)^{1-1/L}\theta_n^{1/(L\rho)}
 +\bar A^L\theta_n^{1/\rho}\right]+n^{-3}.
 \label{eq:localized-population-oracle}
\end{equation}
The fallback is absorbed by the second remainder when its branch
condition fails, after enlarging constants. This theorem is uniform in
all eligible widths/scales; the comparison teacher's scalar
concentration event is common to them.

For the present dictionary, $\rho=\rho_H$ and $v=v_H$. Lemma~\ref{lem:teacher} with the fixed-radius resource~\eqref{eq:fixed-resource} supplies $A=\mathcal A_{m,B,n}$ uniformly for $\lambda\ge1/n$. Its teacher output obeys $\|g\|_\infty^2\le C A^L$: this follows directly from~\eqref{eq:teacher-costs}, since $U_{m,B}^2\le C E_{m,B}^L$.

Choose a population-risk minimizer in the compact comparator parameter box. It exists by continuity and dominated convergence, and has risk at most one because the zero readout is available. Its one scalar Bernstein event, at confidence $\delta/2$, is shared by all learner widths and prior scales. Combine it with the common dictionary event. The preceding bound is exactly~\eqref{eq:oracle}, after absorbing confidence logarithms in $\overline{\mathcal A}$. It controls posterior-averaged raw risk before Jensen, so it also controls the raw mean. This proves both main theorems and the stated sample exponents.


\section{Scope and remaining quantitative gaps}
The result gives a common argument for bounded and unbounded Lipschitz activations. It does not require a margin on every possible input, a robustness restriction on the prior, a prior over smaller architectures, or exact interpolation of noisy labels. It is an agnostic squared-loss theorem with a calibrated Gaussian likelihood.

The sample exponent has quadratic dependence on hidden depth, and the constants need not have polynomial dependence on depth. This is not a polynomial-in-depth theorem. The resulting width and coefficient polynomials are also far from sharp. ReLU homogeneity can lower the comparator energy, while stronger Gaussian scales admit much sharper capacity arguments. The present theorem's purpose is to cover the entire mean-field-or-stronger regime with one proof, including arbitrary globally Lipschitz activations and nonzero $\phi(0)$.

\appendix
\section{A likelihood schedule fixed before seeing the data}
\label{sec:fixed-calibration}
Data dependence is useful for the risk-adaptive terms, but is not essential for a polynomial oracle guarantee.

\begin{corollary}[Deterministic calibration]
\label{cor:deterministic}
Use the same constants as the main theorems. On their root branch, replace the data-dependent solution by the predetermined value
\[
 \lambda_{\rm det}
   =\bigl(\theta_n U_0^{v_H-\rho_H}\bigr)^{1/v_H};
\]
on the fallback branch retain $\lambda_{\rm det}=\omega M$.
This schedule depends on the sample count, confidence level, architecture and known prior, but not on the observations or comparator.

On the nonfallback branch, with probability at least $1-\delta$,
simultaneously for every $N\ge m$ and scale~\eqref{eq:scales},
\begin{align*}
 R(\bar f_{\lambda_{\rm det}})
 \le R_*+(\mathcal A_{m,B,n}+B_d)\lambda_{\rm det}+n^{-3}
 &+C(U_{m,B}+1)\sqrt{\frac{\log(4/\delta)}n}\\
 &+C(U_{m,B}+1)^2\frac{\log(4/\delta)}n.
\end{align*}
The same inequality holds for posterior-averaged raw risk.
Consequently this predetermined Gaussian likelihood already gives a
coefficient-one polynomial oracle, with general-risk exponent
$1/[2H(H+1)]$ for $H\ge2$ and $1/5$ for $H=1$, up to logarithms and
polynomial comparator resources.
\end{corollary}
\begin{proof}
The branch condition gives $1/n\le\lambda_{\rm det}\le1/b_n$.
Set $u=2F_{\lambda_{\rm det}}+b_n\lambda_{\rm det}\le U_0$.
The localized proof only needs the budget inequality
\[
 \lambda_{\rm det}^{v_H}
 \ge\theta_n u^{v_H-\rho_H},
\]
rather than equality. This holds by construction. Its prior-tail argument
and entropy absorption are unchanged, so the same certificate
$F_{\lambda_{\rm det}}+B_d\lambda_{\rm det}+n^{-3}$ follows.
Apply the fixed-radius teacher comparison and scalar Bernstein to a
population-risk-minimizing comparator. On the fallback branch the
deterministic constant certificate remains valid, and that branch
requires only polynomially many samples to exit at fixed depth.
\end{proof}


\begin{thebibliography}{9}
\bibitem{tropp}
J. A. Tropp. User-friendly tail bounds for sums of random matrices.
\emph{Foundations of Computational Mathematics} 12 (2012), 389--434.
\bibitem{gaussian}
M. Ledoux. \emph{The Concentration of Measure Phenomenon}.
American Mathematical Society, 2001.
\end{thebibliography}
\end{document}
